\section{Introduction}\label{sec:introduction}

A dimension count predicts how a subvariety of an abelian variety
meets a fixed torsion coset. The abelian Zilber--Pink conjecture asks
for a finiteness statement when the intersection is larger than this
prediction, even though the torsion coset is allowed to vary. We prove
the conjecture over number fields in arbitrary dimension.

\subsection{The finiteness statement}
Fix a number field $K$, an algebraic closure $\bar K$, and an abelian
variety $A/K$. All subvarieties below are closed and irreducible over
$\bar K$. A \emph{special subvariety} of $A_{\bar K}$ is a torsion coset
$B+\tau$, where $B$ is an abelian subvariety and $\tau$ is a torsion
point. For a subvariety $Z$, write $\langle Z\rangle_{\mathrm{sp}}$ for
the smallest special subvariety containing $Z$.

This smallest subvariety exists: choose a special subvariety of least
dimension containing $Z$, and intersect it with any other such
subvariety. Each irreducible component of a nonempty intersection of
torsion cosets is a torsion coset. The component containing $Z$ has the
same dimension as the chosen one, and hence equals it. This also proves
uniqueness and the implication
\begin{equation}\label{eq:closure-monotonicity}
 Z\subseteq Z'\quad\Longrightarrow\quad
 \langle Z\rangle_{\mathrm{sp}}\subseteq
 \langle Z'\rangle_{\mathrm{sp}}.
\end{equation}
Here the assertion about intersections follows by intersecting the
underlying algebraic subgroups: a nonempty intersection is a coset of
their intersection, torsion points lift through surjective homomorphisms of abelian
varieties, and every component of an algebraic subgroup of an abelian
variety contains a torsion point.

Let $X\subseteq A_{\bar K}$ and put $S=\langle X\rangle_{\mathrm{sp}}$.
An irreducible component $Y$ of $X\cap T$, for a special $T\subseteq S$,
is \emph{atypical for $X$ in $S$} if
\begin{equation}\label{eq:atypical}
 \dim Y>\dim X+\dim T-\dim S.
\end{equation}
Maximality of an atypical subvariety always means maximality for
inclusion among the atypical subvarieties of $X$ in $S$.

\begin{theorem}\label{thm:main}
For every number field $K$, every abelian variety $A/K$, and every
irreducible subvariety $X\subseteq A_{\bar K}$, there are only finitely
many maximal atypical subvarieties of $X$ in
$S=\langle X\rangle_{\mathrm{sp}}$. Equivalently, there are finitely many
proper atypical subvarieties $Y_1,\ldots,Y_m\subsetneq X$ such that every
atypical subvariety is contained in one of the $Y_j$.
\end{theorem}

The list can be empty. No bound on the degrees or heights of the
intersection points is assumed, and no simplicity, endomorphism, or
complex multiplication hypothesis is imposed on $A$.
The theorem concerns torsion cosets in abelian varieties; it gives no
effective or uniform bound for the number of maximal atypical
subvarieties.

For a torsion coset $S$ and an integer $r\ge0$, set
\[
 S^{[r]}=\bigcup_{\substack{T\subseteq S\ \mathrm{special}\\
                         \codim_S T\ge r}}T.
\]
We first prove the following non-density theorem, with universal
quantifiers on the ambient abelian variety and its subvarieties.

\begin{theorem}\label{thm:nondensity}
Let $A$ be an abelian variety over $\Qbar$, and let $X\subseteq A$ be an
irreducible subvariety contained in no proper torsion coset. Then
$X\cap A^{[\dim X+1]}$ is not Zariski dense in $X$.
\end{theorem}

The passage from this assertion to Theorem~\ref{thm:main} uses the
optimal-subvariety reduction of Barroero and Dill. Section~\ref{sec:finiteness}
checks its universal hypotheses and gives the final inclusion argument.
In particular, the proof does not stop with non-density for one fixed
$X$.

The same universal non-density theorem also gives a consequence for
translations by finite-rank subgroups. For point sets
$E,\Gamma\subseteq A(\Qbar)$, write
$E+\Gamma=\{e+\gamma:e\in E,\ \gamma\in\Gamma\}$.

\begin{corollary}[Finite-rank translated intersections]
\label{cor:finite-rank-translates}
Let $A$ be an abelian variety over $\Qbar$, let
$\Gamma\subseteq A(\Qbar)$ be a subgroup satisfying
\[
 \dim_{\mathbb Q}\bigl(\Gamma\otimes_{\mathbb Z}\mathbb Q\bigr)<\infty,
\]
and let $X\subseteq A$ be an irreducible closed subvariety of dimension
$d$ contained in no translate of a proper algebraic subgroup of $A$.
Then
\[
 X(\Qbar)\cap\bigl(A^{[d+1]}(\Qbar)+\Gamma\bigr)
\]
is not Zariski dense in $X$. If $d=1$, this intersection is finite.
\end{corollary}

This is the abelian algebraic-point case of Pink's
Conjecture~5.2 \cite{Pink2005}. The group $\Gamma$ need not be finitely
generated: torsion subgroups and division hulls of finitely generated
subgroups are included. Section~\ref{sec:finite-rank} applies Pink's
reduction \cite[Theorem~5.3]{Pink2005} and records where the stronger
no-proper-translate hypothesis is used.

\subsection{Historical context}
Zilber formulated the unlikely-intersection principle for semiabelian
varieties \cite{Zilber2002}. Pink's semiabelian formulation
\cite[Conjecture~5.1]{Pink2005} specializes to the high-codimension
non-density statement considered here.

The torsion-point case belongs to the Manin--Mumford theorem of Raynaud
\cite[p.~327, Th\'eor\`eme]{Raynaud1983Torsion}: a subvariety with dense
torsion points is a torsion coset. Unlikely intersections allow the torsion points to be
replaced by torsion cosets of varying positive dimensions. Even when
these cosets have large codimension, their intersections can contain
positive-dimensional components; the maximality condition in
Theorem~\ref{thm:main} accounts for this possibility.

For curves in powers of a CM elliptic curve, Viada proved
codimension-two finiteness under the hypothesis that the curve lies in
no translate of a proper algebraic subgroup
\cite[Theorem~2]{Viada2003}. R\'emond and Viada removed this
transversality restriction, retaining the necessary exclusion of proper
algebraic subgroups \cite[Theorem~1.7]{RemondViada2003}. Habegger and Pila proved
finiteness for curves in arbitrary abelian varieties over number fields
\cite[Theorem~1.1]{HabeggerPila2016}. Barroero, K\"uhne and Schmidt
subsequently established the corresponding curve theorem for general
semiabelian varieties over number fields
\cite[Theorem~1.1]{BarroeroKuhneSchmidt2023}.
More recently, Dogra and Pandit obtained bounded-degree results for smooth
curves whose Jacobians surject onto the ambient abelian variety. Together
with established height bounds, these give finiteness. Their local theorem
gives explicit reduction bounds under good-reduction and genus hypotheses,
and cardinality bounds under further restrictions on the simple isogeny factors
\cite[Theorem~1 and Corollary~1]{DograPandit2026}.

Higher-dimensional non-density was known under additional geometric
or height restrictions. For positive-dimensional $X$, consider the
condition
\[
 \dim\pi(X)=\min\{\dim X,\dim(A/B)\}
 \qquad\text{for every quotient }\pi:A\longrightarrow A/B.
\]
In the CM case, R\'emond proved finite-rank translated non-density under
this condition \cite[Theorem~1.3]{Remond2009}; Habegger independently
proved the untranslated case \cite[Corollary~2]{Habegger2009}.
Viada obtained a non-CM elliptic-power case
\cite[Theorem~1.3]{Viada2009}. Habegger and Pila established
non-density under the same quotient condition for arbitrary abelian
varieties \cite[Theorem~1.3(ii)]{HabeggerPila2016}. Without that condition,
their Theorem~1.3(i) gives non-density at every fixed height bound when
$X$ is contained in no proper torsion coset.

There were also higher-dimensional finiteness results for maximal
atypical subvarieties. Habegger and Pila's results imply finiteness of
those of codimension one in $X$
\cite[Theorem~9.14(i)]{HabeggerPila2016}. For codimension-two
subvarieties of powers of an elliptic curve contained in no proper
torsion coset, maximal-atypical finiteness was proved in the CM case by
Checcoli, Veneziano and Viada \cite[Corollary~1.5]{CheccoliVenezianoViada2014}
and without the CM restriction by Hubschmid and Viada
\cite[Theorem~1.1]{HubschmidViada2019}.
Theorem~\ref{thm:main} treats arbitrary abelian varieties and all
codimensions, without the additional quotient-dimension or height
restrictions.

The geometric inputs are Ax's analytic-subgroup theorem
\cite{Ax1972}, in the precise analytic-subset form stated in
\cite[Theorem~5.3]{HabeggerPila2016}, and the finiteness of directions of
geodesic-optimal subvarieties of a fixed parent
\cite[Proposition~6.1]{HabeggerPila2016}.
Habegger and Pila explicitly attribute the latter finiteness to earlier
results of R\'emond \cite[Lemma~2.6 and Proposition~3.2]{Remond2009}.

The height argument follows Vojta's curve method \cite{Vojta1991},
Faltings' extension to abelian varieties \cite{Faltings1991}, and
R\'emond's generalized inequality \cite{Remond2005Vojta}. We use Dill's
version with separate height thresholds for the factors; his introduction
credits Ange with the preceding extension to different projective
varieties \cite{Dill2020}.
The bounded-scale argument adapts the two-measure comparison in the
Bogomolov proofs of Ullmo and Zhang \cite{Ullmo1998,Zhang1998}, based
on Szpiro--Ullmo--Zhang equidistribution \cite{SzpiroUllmoZhang1997}.
Zhang's account \cite[Section~3]{Zhang1998ICM} explicitly traces the
consecutive-difference map to Faltings. The metric-variation calculation
uses Yuan's arithmetic bigness method \cite{Yuan2008}, in the nef adelic
form of Yuan and Zhang \cite{YuanZhang2021}. The local weighted-index
argument follows Faltings' method, and the final cut uses his product
theorem in Evertse's explicit form \cite{Evertse1995Product}.

The estimates proved here keep these inputs uniform across
several height scales and across auxiliary varieties with bounded
fields of definition. The argument for the bounded scale compares
measures only on the fixed low factor; it does not invoke an
equidistribution theorem for a fixed variety as a substitute for
uniformity.

\subsection{Proof strategy and organization}

The proof studies a hypothetical Zariski-generic sequence of
high-codimension intersection points. The square roots of their
N\'eron--Tate heights need not be comparable. Repeatedly taking a quotient
of largest dimension on which the height is smaller separates the ambient variety,
up to isogeny, into at most one bounded level and finitely many successively
larger unbounded levels. Normalizing each level separately makes the
point vectors bounded. Their algebraic subgroup relations then converge
to a real-endomorphism projector that sends the normalized vectors to
zero in the limit.

The auxiliary sequence theorem in Section~\ref{sec:sequences} excludes
the high-codimension limit just obtained. Its geometric input is uniform tangent
positivity, obtained in Section~\ref{sec:tangent} from Ax's theorem and
the finiteness of geodesic-optimal directions in a fixed parent variety.
Projection to smaller abelian varieties handles every possible failure
of this positivity.

There are two arithmetic parts. When every height level is unbounded,
Section~\ref{sec:high} applies Dill's generalized Vojta--R\'emond
inequality. The relevant multiplicative constant is independent of the
accuracy of a rational approximation to the limiting projector.
This permits the approximation error to be made small before taking
points of sufficiently large height.

A bounded level requires a different argument. Section~\ref{sec:minimal}
chooses auxiliary products of minimal dimension whose degrees and fields
of definition are bounded. These bounds concern the varieties, while
the marked points may have arbitrarily large degree. A zero-dimensional
bounded factor reduces to the unbounded case. For a positive-dimensional
bounded factor, Section~\ref{sec:arithmetic} proves a uniform arithmetic
intersection lower bound. The proof compares two limiting measures on
that fixed factor: one from a product map and one from a consecutive
difference map. The first has positive density at a suitable diagonal
point; the second has zero density there. A variational argument on the
changing varieties converts this disagreement into arithmetic positivity.

Small sections supplied by that positivity have a uniformly positive
weighted vanishing index at the marked points
(Section~\ref{sec:index}). The product theorem then yields a proper
factor through a marked point. A multiplicity calculation bounds both
its degree and the number of its Galois conjugates
(Section~\ref{sec:product-cut}). This contradicts minimality and completes the sequence theorem.

\subsection{Conventions}
We identify $\bar K$ with $\Qbar$ and fix an embedding $\Qbar\hookrightarrow
\mathbb C$. A coset without the adjective special may be translated by
an arbitrary point. Its direction is the abelian subvariety being
translated. All logarithmic heights and arithmetic intersection numbers
are absolute, normalized by the degree of the field of definition.
Constants in asymptotic estimates may depend on fixed varieties,
embeddings and polarizations, but not on the sequence index unless
explicitly stated. The notation $u_i\asymp v_i$ means that both ratios
are bounded above by positive constants.
